\chapter{Estado del arte} % Main chapter title
\label{sec:EstadoArte} % For referencing the chapter elsewhere, use
\section{Modelos Basados en Agentes de evasión tributaria}
\label{sec:edo-ABM}

Los Modelos Basados en Agentes aplicados a la evasión tributaria han evolucionado desde formulaciones simples con agentes heterogéneos hacia arquitecturas progresivamente más completas. \textcite{hokamp2010income} introdujeron una sociedad de agentes heterogéneos que incluye ``agentes imitadores tipo b'' (imitating b-type) que observan a sus vecinos en una red social de tamaño $\nu$ y copian las estrategias de evasión exitosas, lo que genera una dinámica social en la que el fraude se propaga por imitación del éxito. Esta arquitectura constituye una base canónica para los modelos posteriores que han integrado el aprendizaje social, incluyendo el modelo de \textcite{canessa2025tax}, que extiende la representación con aprendizaje por refuerzo sobre una vecindad Moore.

Sobre esa base, se han desarrollado marcos computacionales más completos. El modelo SIMULFIS \parencite{llacer2018simulfis, noguera2014tax} integra formalmente tres memorias locales en el agente: memoria de fraude previo, memoria de auditorías previas y memoria de sanciones previas, e incorpora la influencia social mediante una red de agentes enlazados. La regla de actualización combina conformidad social y aprendizaje, lo que permite reproducir patrones agregados de cumplimiento condicionados por la estructura de la red subyacente. El marco de \textcite{hashimzade2014agent} formaliza la elección ocupacional y la interacción entre agentes, mostrando que la topología de la red y el régimen de auditoría interactúan de manera no trivial sobre el ingreso agregado.

La línea evolutiva más reciente formaliza la evasión como juegos evolutivos en redes. \textcite{chica2021evolutionary} modelan la dinámica del fraude en impuestos al consumo mediante una regla de actualización por imitación social en redes power-law, donde cada agente compara su aptitud (fitness) con la de sus vecinos directos y adopta estrategias en función de esa comparación. Esta formulación encaja conceptualmente con estudios previos sobre efectos de vecindad \parencite{arslan2013effects}, en los que individuos recién creados adoptan el comportamiento dominante entre sus vecinos inmediatos (Von Neumann o Moore), lo que acopla la dinámica individual a la estructura espacial subyacente. El papel de la topología se confirma en \textcite{andrei2014agent}, quienes muestran que redes con alta centralidad de cercanía (Power Law y Erdős--Rényi) amplifican el impacto de agentes muy honestos o muy deshonestos, y que auditar a los nodos más conectados (estrategia ``Big Fish'') reduce paradójicamente la recaudación de sanciones porque los hubs propagan normas de cumplimiento que vuelven a la población más honesto.

\section{Memoria social y aprendizaje en red}
\label{sec:edo-memsoc}

La memoria social de los contribuyentes se modela como el número de períodos fiscales recientes cuyos resultados de auditoría (propios o de los vecinos) el agente integra al actualizar su comportamiento. La parametrización más directa de esta memoria es el factor \textit{period-frauds} (PF) introducido por \textcite{canessa2025tax}, que define cuántos períodos hacia atrás observa cada agente a sus vecinos en una vecindad Moore. Su estudio contrasta dos regímenes, PF$=1$ (memoria corta) y PF$=5$ (memoria larga), y muestra que la extensión de la memoria social modula críticamente la dinámica agregada: con memoria corta, los agentes reaccionan sobre todo a señales muy recientes y la estrategia de selección de auditoría resulta determinante; con memoria larga, en cambio, la información histórica se acumula en las creencias subjetivas de los agentes, lo que tiende a estabilizar normas de cumplimiento y reduce la ventaja comparativa de las estrategias de focalización. Este hallazgo es el avance crítico sobre el que se construye la presente tesis.

La memoria social puede también entenderse como un mecanismo de contagio. \textcite{dimauro2019contagion} modelan la evasión tributaria como un juego de contagio en una red de mundo pequeño y describen un umbral crítico: cuando la fracción inicial de contribuyentes cumplidores supera aproximadamente 55\,\%, el sistema converge hacia un equilibrio de alto cumplimiento; por debajo de ese valor, la evasión se vuelve dominante. El umbral depende de dos parámetros individuales: el ``Imitation Factor'' (sensibilidad al comportamiento vecinal) y el ``Capital Factor'' (sensibilidad a las condiciones económicas personales), lo que confirma que la memoria social y la experiencia individual interactúan de forma no lineal.

La evidencia experimental respalda parcialmente estos mecanismos. \textcite{alm2017neighbour} muestran que información sobre vecinos tiene efectos estadísticamente significativos y económicamente grandes en el cumplimiento, pero la dirección del efecto depende del contenido de la información: informar sobre los resultados de auditorías de vecinos (sanciones pagadas) aumenta tasas de declaración y reporte, mientras que informar sólo sobre tasas de declaración de los vecinos puede reducirlas. La literatura inserta así un matiz importante: la memoria social no es uniformemente pro-cumplimiento y su efecto depende de qué información se propague.

Finalmente, hay modelos que hacen explícita la mediación entre memoria individual y memoria social. \textcite{garcia2019network} proponen que cada agente actualiza su probabilidad subjetiva de auditoría promediando tres canales: su tasa subjetiva previa (prior), el recuerdo de sus propias auditorías pasadas (memoria individual) y las señales recibidas de sus vecinos (influencia social). En este modelo los agentes comunican a sus vecinos sólo su memoria sobre auditorías previas y no información sobre esquemas de evasión completos, lo que ofrece un matiz sobre el alcance de la información socialmente compartida.

\section{Estrategias de selección de auditoría}
\label{sec:edo-estrategias}

La selección de contribuyentes a auditar constituye una decisión de política y ha sido objeto de modelación comparativa. \textcite{hashimzade2014agent} contrastan cuatro estrategias en un ABM en el que los agentes eligen entre empleo asalariado y autoempleo: auditar a todos con la misma probabilidad, auditar a un número fijo de contribuyentes por ocupación en cada período, y dos variantes en las que la autoridad reparte sus auditorías entre ocupaciones de forma alternada o aleatoria. Los autores encuentran que mantener un número fijo de auditorías por ocupación en cada período rinde más que las demás estrategias, porque la caída del recaudo en los períodos con poca auditoría no se compensa con los períodos de auditoría intensa.

El uso de análisis predictivos para focalizar auditorías ha sido formalizado por \textcite{hashimzade2016predictive}, que muestran que el empleo de regresión predictiva sobre datos fiscales genera un aumento significativo en el recaudo respecto a la auditoría aleatoria, mediante una regla efectiva cercana a un punto de corte basado en ingreso declarado: la casi totalidad del esfuerzo de auditoría se concentra en los percentiles más bajos de declaración. La extensión conductual de \textcite{hashimzade2016risk} adopta un modelo de regresión Tobit y asume un ``efecto objetivo'' (target effect), según el cual el contribuyente auditable eleva su probabilidad subjetiva de auditoría, mientras que la evidencia sobre el alternativo ``efecto cráter de bomba'' (bomb-crater effect, en el que la auditoría reduce el riesgo percibido) sigue siendo mixta.

La introducción de un tercer tipo de estrategia, basada en puntajes de riesgo computados secuencialmente, es el núcleo del modelo de \textcite{canessa2025tax}. Sus autores contrastan tres enfoques: auditoría aleatoria, focalización por nivel de ingreso y focalización estratégica por puntaje de riesgo, encontrando que el desempeño relativo de cada estrategia depende de la extensión de la memoria social de los contribuyentes: con memoria social corta, la focalización estratégica supera a la aleatoria, mientras que con memoria social larga la auditoría aleatoria alcanza desempeño marginalmente equivalente y la focalización por nivel de ingreso rinde de forma consistentemente inferior. Profundizar en esta interacción es precisamente lo que motiva la presente tesis.

\section{Memoria institucional y revisión retrospectiva}
\label{sec:edo-meminst}

La literatura reconoce de manera genérica que la memoria histórica de sanciones previas es crucial para sostener la disuasión \parencite{antunes2005tax, bloomquist2004multi, hashimzade2018using}, pero la modelación explícita de la revisión retrospectiva del auditor como variable independiente es reciente y fragmentaria. El estudio más directamente relevante es el de \textcite{hokamp2014dynamics}, que formalizan el back auditing con el parámetro $\alpha$ (el número de períodos previos sujetos a auditoría) y muestran que el back auditing domina la dinámica de evasión en todos los escenarios simulados: sin retroauditoría ($\alpha = 0$) el fraude permanece elevado y volátil, mientras que la retroauditoría reduce drásticamente la evasión. Verificado contra el manuscrito fuente, este resultado se sostiene en el plano cualitativo, mientras que las cifras específicas sobre magnitud de reducción son atribuibles a la simulación de referencia de \textcite{hokamp2010income}.

En el caso chileno, el Servicio de Impuestos Internos puede revisar declaraciones de varios períodos fiscales anteriores al declarado, en una práctica institucional que formaliza la profundidad retrospectiva del fiscalizador. El Plan de Gestión de Cumplimiento Tributario~\parencite{sii2026pgct} recoge esta facultad como parte de la estrategia de control y describe la necesidad de procesar historiales de comportamiento para sostener un alto impacto disuasivo con el menor costo administrativo posible. Esta práctica define el rango de variación que la presente tesis adopta como dominio experimental del parámetro de profundidad.

El ruido de detección sobre información histórica envejecida es un componente empíricamente documentado de la memoria institucional. \textcite{debacker2018once} explotan auditorías aleatorias del Internal Revenue Service (2006--2009) y muestran que el efecto disuasivo de una auditoría se atenúa a medida que la información histórica envejece, con una caída más pronunciada para ingresos sin reporte de terceros y más gradual para ingresos salariales. La volatilidad del ingreso del contribuyente acelera este decaimiento. En conjunto, los resultados describen un patrón cuantitativo que justifica modelar la degradación de la información como un mecanismo explícito de la memoria institucional.

La modelación teórica reciente ha comenzado a integrar explícitamente la noción de ``anclas de auditoría'' (audit anchors), entendidas como registros institucionales previos que guían selectivamente las decisiones de fiscalización, aunque su tratamiento sistemático en ABM permanece limitado.

\section{Síntesis del vacío}
\label{sec:edo-vacio}

Esta revisión permite identificar un vacío rodeado por tres frentes convergentes. Primero, los modelos de memoria social demuestran que el comportamiento de los contribuyentes se moldea por la observación de sus vecinos, pero su efecto agregado depende del contenido y la persistencia temporal de la información \parencite{andrei2014agent, dimauro2019contagion, alm2017neighbour, garcia2019network}. Segundo, las estrategias de fiscalización han sido comparadas bajo distintos supuestos de aprendizaje, y su eficiencia relativa está mediatizada por la memoria social \parencite{hashimzade2014agent, hashimzade2016predictive, canessa2025tax}. Tercero, la memoria institucional ha sido tratada de forma fragmentaria: el back auditing domina cualitativamente \parencite{hokamp2014dynamics} y su efecto decae con la antigüedad \parencite{debacker2018once}, sin que la literatura haya integrado en un mismo ABM la interacción de la profundidad retrospectiva con el ruido de detección. El estudio más cercano al vacío es el de \textcite{canessa2025tax}, que caracteriza la interacción entre memoria social y estrategia de fiscalización pero mantiene un fiscalizador puramente transaccional, sin parametrizar la profundidad retrospectiva del auditor. En particular, se desconoce bajo qué configuraciones de memoria institucional y memoria social cada estrategia de auditoría (aleatoria, focalizada por nivel de ingreso o por riesgo) resulta más efectiva en términos de cumplimiento tributario agregado a largo plazo.
